\section{Introducción}\label{sec:intro}

\subsection{Contexto}

En el ámbito de la medicina predictiva, la efectividad de los modelos de Machine Learning (ML) supervisados no reside únicamente en la sofisticación de sus arquitecturas algorítmicas, sino en la integridad y el tratamiento metodológico de los datos que los alimentan.

Los datasets clínicos y biomédicos presentan desafíos intrínsecos que los
diferencian de otros dominios: altos volúmenes de valores faltantes (missing
values), distribuciones asimétricas, multicolinealidad entre biomarcadores y,
fundamentalmente, el riesgo de Data Leakage. En medicina, un
dato faltante rara vez es un error de registro aleatorio; a menudo, es una señal
clínica o un protocolo hospitalario específico.
Por tanto, el preprocesamiento se convierte en la etapa más
crítica del ciclo de vida del dato, pues una decisión puramente técnica puede
destruir la coherencia biológica necesaria para que un modelo sea útil en un
entorno clínico real.

En la construcción de flujos de trabajo (pipelines) de ciencia de datos, es habitual el empleo de estrategias de preprocesamiento de naturaleza técnica, tales como la imputación basada en tendencias centrales globales (medias o medianas), la eliminación de variables según umbrales estadísticos de nulidad o la selección de características mediante algoritmos de aprendizaje. Si bien estos métodos aseguran la eficiencia computacional, plantean un escenario de interés para la investigación: evaluar el impacto diferencial que aporta la integración del razonamiento clínico-estadístico en la preparación de los datos.

Por lo tanto, esta investigación se centra en la necesidad de establecer un marco comparativo robusto o benchmarking. El objetivo es cuantificar y analizar cómo influye un enfoque de preprocesamiento informado por la lógica médica y el análisis contextual del dataset en el rendimiento, la estabilidad y la fiabilidad de los modelos de aprendizaje automático. De este modo, el estudio busca determinar si esta personalización del pipeline genera diferencias significativas en la capacidad de discriminación diagnóstica frente a las estrategias técnicas estandarizadas aplicadas a variables clínicas binarias.

\subsection{Hipótesis}\label{sec:hipotesis}

Sobre la base del problema expuesto, se plantea la siguiente hipótesis de
investigación:

Las estrategias de preprocesamiento diseñadas mediante el uso de razonamiento
clínico-estadístico y un análisis contextual profundo del dataset generan
resultados predictivos superiores y modelos más robustos en comparación con las
estrategias puramente técnicas o automáticas aplicadas sobre datasets biomédicos
complejos.

Se asume que la integración del conocimiento experto permite identificar
variables informativas ocultas en el ruido y mitigar sesgos metodológicos que
las reglas estadísticas estándar no son capaces de detectar.

\subsection{Objetivos}

\textbf{Objetivo general}

Comparar el impacto de diferentes estrategias de preprocesamiento de datos sobre
el rendimiento de modelos de aprendizaje supervisado aplicados a la predicción
de variables clínicas binarias en un entorno de cuidados críticos.

\textbf{Objetivos específicos}

\begin{itemize}
  \item Diseñar una estrategia clínico-estadística (V1): Construir un flujo de
    procesamiento fundamentado en la lógica médica, utilizando imputaciones
    estratificadas por diagnóstico y una depuración selectiva de variables
    basada en la prevención de Data Leakage y la coherencia biológica.
  \item Diseñar una estrategia técnico-estadística (V2): Desarrollar un flujo de
    procesamiento basado en heurísticas técnicas convencionales y criterios de
    tendencia central global, actuando como grupo de control agnóstico al
    dominio médico.
  \item Construir e implementar un sistema de benchmarking: Desarrollar una
    infraestructura automatizada en R/tidymodels que permita el entrenamiento y
    evaluación sistemática de diversos algoritmos (Random Forest, GBM, Logistic
    Regression, entre otros) sobre las distintas versiones del dataset.
  \item Comparar y evaluar el rendimiento de los modelos: Analizar la capacidad de
    discriminación de los modelos mediante métricas de área bajo la curva
    (AUC-ROC), sensibilidad y especificidad en cada escenario de
    preprocesamiento.
  \item Analizar el impacto del preprocesamiento: Determinar cuantitativa y
    cualitativamente en qué medida el razonamiento clínico influye en la
    estabilidad y precisión de las predicciones frente a las estrategias
    automáticas.
\end{itemize}
